\chapter{Introdução}
\label{chap:introducao}

A administração de sistemas e redes vive hoje uma transformação profunda, impulsionada por exigências cada vez maiores de disponibilidade, escalabilidade e rastreabilidade. Gerir infraestruturas à mão, serviço a serviço, tornou-se um risco operacional que muitas organizações já não podem assumir. A necessidade de implementar plataformas de forma rápida, reprodutível e resistente a falhas humanas tem forçado a transição para tecnologias de \textit{containerização} e para a filosofia de \textit{Infrastructure as Code} (IaC) \cite{morris2016infrastructure}.

\section{Contextualização e Motivação}

Em infraestruturas de pequena ou média dimensão --- um cenário frequente em laboratórios de investigação, PMEs ou ambientes de ensino ---, a gestão manual de serviços críticos como DNS, \textit{reverse proxies}, VPNs ou servidores de ficheiros torna-se rapidamente insustentável. Cada intervenção direta no sistema é um ponto de falha: ficheiros de configuração modificados sem rasto, dependências instaladas que não figuram em documentação ou serviços repostos com parâmetros divergentes dos originais são erros comuns em ambientes que carecem de automação.

A motivação deste projeto surgiu precisamente desta realidade. Durante a implementação, a ausência de uma base de automação agravou os tempos de recuperação após incidentes. Um exemplo claro foi a simples reinicialização do \textit{router} da rede local, que derrubou a conectividade da máquina virtual (VM) e exigiu uma intervenção manual de diagnóstico e reconfiguração. Num ambiente automatizado, esta intervenção seria evitável ou resolvida em instantes. Este incidente, que será analisado detalhadamente no Capítulo \ref{chap:implementacao}, serve como prova da relevância prática dos objetivos deste trabalho.

O paradigma de \textit{containerização}, através de ferramentas como o Docker, resolve estes problemas ao encapsular aplicações e serviços em unidades portáteis e isoladas. As dependências são declaradas de forma explícita, assegurando que o ambiente de execução é consistente, independentemente da máquina física hospedeira \cite{docker_docs}. Quando combinamos esta portabilidade com a automação via Ansible e \textit{pipelines} de CI/CD, garantimos que qualquer alteração à infraestrutura é auditada, testada e aplicada de forma controlada, evitando o "desvio de configuração" (\textit{configuration drift}) \cite{humble2010continuous}.

\section{Objetivos}

O objetivo central deste trabalho é construir, configurar e documentar uma plataforma de serviços de rede totalmente baseada em \textit{containers}, orquestrada com Docker Compose, focada em segurança, observabilidade e CI/CD.

De forma específica, os objetivos operacionais definidos são:

\begin{itemize}
    \item \textbf{Implementar serviços de rede em \textit{containers}:} Servidor DNS (Pi-hole), \textit{Reverse Proxy} HTTPS (Nginx), servidor Web (Nginx), VPN (WireGuard), Servidor de E-mail (Docker-Mailserver), servidor de ficheiros (Nextcloud) e \textit{stack} de monitorização (Prometheus, Grafana, Loki).
    
    \item \textbf{Garantir reprodutibilidade total:} A infraestrutura deve ser recriável a partir do repositório Git sem intervenção manual, seguindo o princípio da \textit{Single Source of Truth} \cite{morris2016infrastructure}.
    
    \item \textbf{Aplicar práticas rigorosas de segurança:} Utilização de volumes \textit{read-only}, segmentação de redes internas Docker, encriptação de segredos com \textit{Ansible Vault} e reforço de tráfego via TLS/WireGuard.
    
    \item \textbf{Automatizar o \textit{deploy} com CI/CD:} Qualquer \textit{commit} na \textit{branch} principal deve despoletar automaticamente o \textit{deploy} da infraestrutura, assegurando consistência entre o código e o ambiente de produção.
    
    \item \textbf{Integrar observabilidade centralizada:} Implementação de recolha de métricas de \textit{hardware} e logs centralizados, com visualização em painéis Grafana para monitorização proativa.
    
    \item \textbf{Documentar procedimentos de \textit{troubleshooting}:} Registo detalhado dos desafios técnicos e erros encontrados durante o projeto, demonstrando o raciocínio aplicado na sua resolução.
\end{itemize}

\section{Metodologia}

A metodologia adotada divide-se em três eixos principais. Primeiro, uma estratégia incremental: os serviços foram implementados por camadas, começando pela infraestrutura base (SO e Docker), progredindo para os serviços de rede (DNS, Web, VPN, E-mail) e culminando na observabilidade e automação. Esta estratificação permitiu validar cada camada antes de avançar para a próxima.

Segundo, o versionamento contínuo: toda a configuração foi mantida em Git, permitindo o rastreio de alterações e a reversão imediata de configurações problemáticas. Por fim, uma documentação ativa: cada obstáculo técnico foi registado no momento da ocorrência, preservando o contexto necessário para que este relatório sirva não apenas como descrição, mas como \textit{runbook} de referência para futuras intervenções.

\section{Estrutura do Documento}

O relatório organiza-se da seguinte forma:

\begin{itemize}
    \item \textbf{Capítulo \ref{chap:estado_arte} --- Estado da Arte:} análise das tecnologias escolhidas face às alternativas de mercado.
    \item \textbf{Capítulo \ref{chap:necessidades} --- Levantamento das Necessidades:} requisitos funcionais e análise da arquitetura base.
    \item \textbf{Capítulo \ref{chap:implementacao} --- Arquitetura e Implementação:} detalhe técnico da montagem da infraestrutura e resolução de problemas.
    \item \textbf{Capítulo \ref{cap:test} --- Testes e Avaliação:} validação prática dos serviços e discussão dos resultados.
    \item \textbf{Capítulo \ref{cap:conclusions} --- Conclusões:} síntese do trabalho e propostas de evolução futura.
\end{itemize}